\subsubsection{Rovinné grafy}
\label{sssec:rovinne-grafy}

Přívlastek \uv{rovinný} jsme v kontextu kreslení grafů zmínili již několikrát a
neformálně jsme uvedli, že grafu řekneme rovinný, když se dá nakreslit bez
křížení hran. Formulace \uv{dá nakreslit} je zde zásadního významu. Libovolný
graf, který má aspoň dvě hrany, lze vždy nakreslit tak, aby se tyto křížily. Nás
však zajímá pouze, zda existuje \textbf{nějaké} nakreslení takového grafu, ve
kterém se žádné hrany nekříží.

Obecně, minimální nutný počet křížení hran v nakreslení grafu je bohatě nejvíce
studovaná vlastnost z celé oblasti kreslení grafů. Definice křížení (nebo
absence téhož) v nakreslení grafu není výrazně odlišná od té intuitivní.
Zkrátka, řekneme, že dvě hrany se nekříží, když v průniku jejich obrazů leží
pouze nakreslené vrcholy (tam se samozřejmě hrany křížit mohou a některé
musejí).

\begin{definition}[Rovinný graf]
 \label{def:rovinny-graf}
 Graf $G = (V,E)$ nazveme \emph{rovinným}, když existuje jeho nakreslení $D =
 (p,c)$ takové, že pro každé dvě různé hrany $e,e' \in E$ platí
 \[
  \img \gamma_e \cap \img \gamma_{e'} \subseteq \img p.
 \]
 Každé nakreslení rovinného grafu, které splňuje tuto podmínku, nazveme též
 \emph{rovinným}.
\end{definition}

Ve zbytku kapitoly nás čekají jeden snadný a dva těžké úkoly.
\begin{enumerate}
 \item Definovat několik základních operací na vrcholech a hranách grafu, které
  přijdou vhod v dalším textu.
 \item Zformulovat a dokázat několik základních vlastností rovinných gra\-fů.
 \item Dokázat Kuratowskiho větu, která dokonale klasifikuje rovinné gra\-fy,
  tj. poskytuje konkrétní přímočaré kritérium, podle nějž lze poznat, zda je
  graf rovinný, či nikoliv.
\end{enumerate}

Začněme bodem (1). Budeme potřebovat zavést operace přidání a odebrání vrcholů a
hran a konečně operaci dělení hrany. Význam prvních čtyř je snad jasný z názvu,
pátá operace způsobuje vlastně přidání vrcholu doprostřed nějaké hrany. Dá se
pomocí prvních čtyř vyjádřit jako
\begin{itemize}
 \item odebrání dlouhé hrany,
 \item přidání vrcholu,
 \item přidání dvou hran vedoucích z tohoto vrcholu do počátečního a koncového
  vrcholu původní odebrané hrany.
\end{itemize}

I když intuitivně se operace pojímají snadno, formální definice není naprosto
přímočará. Spěšně si ji sepíšeme.

\begin{definition}[Základní grafové operace]
 \label{def:zakladni-grafove-operace}
 Ať $G = (V,E)$ je graf. Definujeme následující operace na množinách $V$ a $E$:
 \begin{itemize}
  \item operace přidání vrcholu, zapisujeme
   \[
    G + v \coloneqq (V \cup \{v\}, E).
   \]
  \item operace odebrání vrcholu (zde je třeba \textbf{odebrat i všechny hrany,
   které do tohoto vrcholu vedou}), pro $v \in V$ zapisujeme
   \[
    G - v \coloneqq (V \setminus \{v\}, E \setminus \{e \in E \mid v \in e\}).
   \]
  \item operace přidání hrany, pro $u,v \in V$ takové, že $uv \notin E$,
   zapisujeme
   \[
    G + uv \coloneqq (V, E \cup \{uv\}).
   \]
  \item operace odebrání hrany, pro $uv \in E$ zapisujeme
   \[
    G - uv \coloneqq (V, E \setminus \{uv\}).
   \]
  \item operace dělení hrany, pro $uv \in E$ zapisujeme
   \begin{equation*}
    \begin{split} 
     G \de uv &\coloneqq (V \cup \{w\}, (E \setminus \{uv\}) \cup \{uw\} \cup
     \{wv\})\\
              &= (((G - uv) + w) + uw) + wv.
    \end{split}
   \end{equation*}
 \end{itemize}
\end{definition}

\begin{figure}[h]
 \centering
 \begin{tikzpicture}[scale=1.5]
  \node[vertex] (u) at (0,0) {};
  \node[vertex] (v) at (2,-1) {};
  \node[left=0mm of u] {$u$};
  \node[right=0mm of v] {$v$};
  \draw[thick] (u) -- (v);
  
  \draw[thick,->] (3,-0.5) to node[midway,yshift=4mm] {$G \de uv$} (4,-0.5);
  
  \node[vertex] (u2) at (5,0) {};
  \node[vertex] (v2) at (7,-1) {};
  \draw[thick,myred] (u2) to node[midway,myred,vertex] (w) {} (v2);
  \node[left=0mm of u2] {$u$};
  \node[right=0mm of v2] {$v$};
  \node[below left=-1mm and -1mm of w,myred] {$w$};
  
 \end{tikzpicture}

 \caption{Operace dělení hrany $uv$.}
 \label{fig:deleni-hrany}
\end{figure}

Pokračujeme bodem (2). Asi není příliš překvapivé, že rovinné grafy nemohou mít
příliš mnoho hran (vzhledem k počtu vrcholů). Čím víc hran do grafu přidám, tím
se snižuje šance, že každou další hranu zvládnu nakreslit bez křížení s
ostatními.

Nejprve objasníme, k čemu nám vlastně slouží
\hyperref[thm:jordanova-o-kruznici]{Jordanova věta o kružnici}. Totiž, kromě
křivek a bodů reprezentujících hrany a vrcholy, získává nakreslení
\textbf{rovinných} grafů ještě jednu strukturu -- stěny. Lidsky řečeno, stěny
daného nakreslení jsou oblasti roviny ohraničené hranami, přesněji cestami.
Nahlédneme, že v důsledku \hyperref[thm:jordanova-o-kruznici]{Jordanovy věty} se
po vynětí všech hran nakreslení daného grafu $G = (V,E)$ rozpadne rovina na
několik omezených oblastí (tzv. \uv{vnitřní stěny}) a jednu neomezenou (tzv.
\uv{vnější stěna}). Každou z těchto oblastí nazveme \emph{stěnou} grafu $G$ a
množinu všech stěn označíme písmenem $F$ (z angl. \textbf{f}ace).

\begin{figure}[h]
 \centering
 \begin{tikzpicture}
  \pgfdeclarelayer{background}
  \pgfdeclarelayer{foreground}
  \pgfsetlayers{background,main,foreground}
  \begin{pgfonlayer}{foreground}
   \node[vertex] (a) at (0,0) {};
   \node[vertex] (b) at (1,2) {};
   \node[vertex] (c) at (3,3) {};
   \node[vertex] (d) at (2,-1) {};
   \node[vertex] (e) at (4,2) {};
   \node[vertex] (f) at (4,0) {};
   
   \node[vertex] (k) at (-0.5,2.5) {};
   \node[vertex] (l) at (-1,1.5) {};
   \node[vertex] (m) at (5,-0.5) {};
   
   \draw[thick] (a) -- (b);
   \draw[thick] (b) -- (c);
   \draw[thick] (c) to node[midway,vertex] {} (d);
   \draw[thick] (d) -- (a);
   \draw[thick] (c) -- (e);
   \draw[thick] (e) -- (f);
   \draw[thick] (f) -- (d);

   \draw[thick] (b) -- (k);
   \draw[thick] (k) -- (l);
   \draw[thick] (f) -- (m);
   
  \end{pgfonlayer}
  \fill[myred,opacity=0.25] (a.center) -- (b.center) -- (c.center) --
   (d.center) -- (a.center);
  \fill[mygreen,opacity=0.25] (c.center) -- (e.center) -- (f.center) -- (d.center)
   -- (c.center);
  \node[myred] (F1) at (1.5,0.8) {$\mathbf{f_1}$};
  \node[mygreen] (F2) at (3.3,0.8) {$\mathbf{f_2}$};

  \begin{pgfonlayer}{background}
   \fill[myblue,opacity=0.25] (-4,4) -- (8,4) -- (8,-2) -- (-4,-2) -- (-4,4);
   \node[myblue] (F3) at (6,3) {$\mathbf{f_3}$};
  \end{pgfonlayer}
  
 \end{tikzpicture}

 \caption{Nakreslení grafu s vnitřními stěnami $\mathbf{\clr{f_1}}$ a
  $\mathbf{\clg{f_2}}$ a vnější stěnou $\mathbf{\clb{f_3}}$.}
 \label{fig:steny-nakresleni}
\end{figure}

\begin{remark}
 Všimněte si, že pojem \emph{stěny} definujeme pouze pro \textbf{rovinné} grafy.
 To má vlastně dva důvody. Zaprvé, stěny, které mají na svých hranicích
 průsečíky nakreslených hran neodpovídající žádným vrcholům, vytvářejí strukturu
 nezávislou na grafu $G = (V,E)$, jejž kreslíme. To je algebraicky zcela
 nepřirozené.

 Druhý důvod je více praktický. Totiž, jak jsme již zmínili, hrany můžeme
 kreslit tak, aby vznikl libovolný počet křížení. To zároveň znamená, že můžeme
 vytvořit libovolný počet stěn v daném nakreslení. Zásadní výsledek pro rovinné
 grafy (který si dokážeme) říká, že počet stěn v každém jeho rovinném nakreslení
 závisí pouze na počtu vrcholů a hran, a \textbf{nikoli na volbě samotného
 rovinného nakreslení}. To je překvapivě velmi hluboký výsledek, neboť ukazuje,
 že -- v principu geometrická -- struktura stěn je v případě rovinných grafů
 naprosto kompatibilní s jeho -- v principu algebraickou -- strukturou hran.
\end{remark}

Abychom směli tvrdit, že graf vůbec \emph{má} nějaké stěny, potřebujeme si
rozmyslet, že nakreslený cyklus je topologická kružnice. To je intuitivně
zřej\-mé, zkrátka na sebe postupně napojíme nakreslené hrany a ta poslední bude
končit tam, kde ta první začala. Rozmyslet si tento fakt formálně je mírně
složitější úloha.

Potřebujeme definovat spojení dvou křivek $\gamma_1,\gamma_2$ takových, že
$\gamma_1(1) = \gamma_2(0)$ a dokázat, že za předpokladu, že se neprotínají
nikde jinde, se jedná rovněž o křivku v rovině. Princip definice je velmi
přímočarý -- nová křivka $\gamma$ bude zkrátka v čase od $0$ do $1 / 2$ sledovat
křivku $\gamma_1$ a od $1 / 2$ do $1$ sledovat křivku $\gamma_2$. Formálně lze
toto nové zobrazení $\gamma:[0,1] \to \R^2$ definovat například následovně.

\begin{definition}[Spojení křivek]
 \label{def:spojeni-krivek}
 Ať $\gamma_1,\gamma_2$ jsou dvě křivky v rovině, které se neprotínají, a platí
 $\gamma_1(1) = \gamma_2(0)$. Potom definujeme jejich \emph{spojení}, zapisované
 často zkrátka $\gamma_1\gamma_2$, předpisem
 \[
  \gamma_1\gamma_2(t) \coloneqq
  \begin{cases}
   \gamma_1(2t), &\text{pokud } t \in [0,\frac{1}{2}]\\
   \gamma_2(2t - 1), &\text{pokud } t \in [\frac{1}{2},1].
  \end{cases}
 \]
\end{definition}

\begin{lemma}
 \label{lem:spojeni-krivek}
 Ať $\gamma_1,\gamma_2$ jsou křivky v rovině, které se neprotínají, a
 $\gamma_1(1) = \gamma_2(0)$. Potom je spojení $\gamma_1\gamma_2$ dobře
 definované a je to křivka v~rovině.
\end{lemma}
\begin{proof}
 \uv{Dobrá definovanost} $\gamma_1\gamma_2$ zkrátka znamená, že jsme nenapsali
 žádný nesmysl, tj. že to je opravdu \emph{zobrazení} (to je zde zřejmé) a že
 jsme při jeho definici například neuvažovali body mimo domény zobrazení
 $\gamma_1$ a $\gamma_2$. To je velmi snadné ověřit, neboť pro $t \in [0,1 / 2]$
 je $2t \in [0,1]$ a pro $t \in [1 / 2,1]$ je $2t-1 \in [0,1]$. V čase spojení
 $t = 1 / 2$ máme
 \[
  \gamma_1(2t) = \gamma_1(1) = \gamma_2(0) = \gamma_2(2t - 1),
 \]
 tedy vše funguje, jak má.

 Fakt, že $\gamma_1\gamma_2$ je křivka v rovině, je též snadno vidět. Totiž,
 zcela jistě je to zobrazení $[0,1] \to \R^2$. Dále, je spojité, ježto
 $\gamma_1$ i $\gamma_2$ jsou spojitá a plynule na sebe navazují. Je též prosté,
 neboť $\gamma_1$ i $\gamma_2$ jsou prostá a z~předpokladu se neprotínají.
\end{proof}

Z \hyperref[lem:spojeni-krivek]{předchozího lemmatu} triviální indukcí vyplývá,
<<<<<<< Updated upstream
že je-li 
=======
že jsou-li $\gamma_1,\gamma_2,\ldots,\gamma_n$ křivky v rovině takové, že
$\gamma_i(1) = \gamma_{i+1}(0) \; \forall i$ a žádné dvě křivky se vzájemně
neprotínají, pak $\gamma_1\gamma_2\cdots \gamma_n$ je též křivka v~rovině. Tento
fakt již přirozeně vede k využití \hyperref[thm:jordanova-o-kruznici]{Jordanovy
věty} pro rovinná nakreslení. Máme-li totiž posloupnost
$\gamma_1,\ldots,\gamma_n$ jako výše se stejnými podmínkami až na jedinou
výjimku, že $\gamma_n(1) = \gamma_1(0)$, pak je tato topologickou kružnicí.
Všimněme si též, že uvažovaná posloupnost je vždy nakreslením nějakého cyklu.

Nyní již konečně můžeme přikročit k formální definici \emph{stěny} rovinného
nakreslení grafu.

\begin{definition}[Stěna rovinného nakreslení]
 \label{def:stena-rovinneho-nakresleni}
 Ať $G = (V,E)$ je rovinný graf a $D = (p,c)$ nějaké jeho rovinné nakreslení.
 Označme písmenem $\mathcal{C}$ množinu všech cyklů grafu $G$. Každému cyklu
 $e_1\cdots e_n \in \mathcal{C}$ odpovídá topologická kružnice $\gamma_1\cdots
 \gamma_n$, neboť uvažované nakreslení je rovinné, a tedy se $\gamma_i$ navzájem
 neprotínají.

 Ať $\kappa_{C}$ značí topologickou kružnici vzniklou nakreslením cyklu $C \in
 \mathcal{C}$. Dle \hyperref[thm:jordanova-o-kruznici]{Jordanovy věty}
 rozděluje $\kappa_{C}$ rovinu na dvě oblasti. Každou oblast množiny
 \[
  \R^2 \setminus \bigcup_{C \in \mathcal{C}} \img \kappa_C
 \]
 nazýváme \emph{stěnou} nakreslení $D$. Množinu všech stěn daného nakreslení $D$
 značíme $F_D$.
\end{definition}

Vlastnosti rovinných nakreslení grafů je obecně obtížné kvantifikovat.
E\-xistuje vlastně jen jeden obecný vzorec, pomocí kterého se pak dokazuje řada
dalších tvrzení. Ujal se pro něj název \uv{Eulerův vztah}.

\begin{claim}[Eulerův vztah]
 \label{claim:euleruv-vztah}
 Ať $G = (V,E)$ je rovinný graf s $k \in \N$ komponentami souvislosti a $D$
 nějaké jeho rovinné nakreslení. Pak platí
 \[
  \# V - \# E + \# F_D = k + 1.
 \]
\end{claim}
\begin{proof}
 Nejprve budeme předpokládat, že $G$ je souvislý, tj. $k = 1$ a dokážeme, že
 \[
  \# V - \# E + \# F_D = 2.
 \]
 Je-li $G$ navíc strom, pak $\# V = \# E + 1$ a $\# F_D = 1$, ježto $G$
 neobsahuje cyklus. Odtud dostaneme
 \[
  \# V - \# E + \# F_D = \# E + 1 - \# E + 1 = 2.
 \]
 Ať $G$ nyní naopak obsahuje cyklus. Postupujme indukcí podle počtu hran, $\#
 E$. Prvním smysluplným případem je $\# E = 3$, neboť cyklus musí mít aspoň tři
 hrany. Pak $G = K_3$, tj. trojúhelník (úplný graf na třech vrcholech), a vzorec
 platí. Dále ať $e \in E$ je hrana obsažená v nějakém cyklu. Z indukčního
 předpokladu vzorec platí pro graf $G - e$. Nakreslení hrany $e$ vytvoří cyklus,
 a tudíž z \hyperref[thm:jordanova-o-kruznici]{Jordanovy věty} rozdělí
 existující stěnu na dvě nové. Protože počet hran stoupne o 1, počet stěn rovněž
 a počet vrcholů zůstane nezměněn, platí vzorec i pro $G$.

 Ať nyní $G$ není nutně souvislý graf s $k$ komponentami souvislosti značenými
 $G_i = (V_i,E_i)$, $i \leq k$. Zřejmě platí
 \[
  \# V = \sum_{i = 1}^{k} \# V_i \quad \text{a} \quad \# E = \sum_{i=1}^{k} \#
  E_i,
 \]
 neboť komponenty grafu z definice nesdílejí hrany ani vrcholy.

 Označme $D_i$ nakreslení $i$-té komponenty. Každá komponenta přispívá do
 celkového počtu stěn $G$ všemi svými vnitřními (tedy omezenými) stěnami, avšak
 všechny komponenty sdílejí vnější neomezenou stěnu. Sečteme-li tedy velikosti
 $\# F_{D_i}$ přes všechna $i$, pak onu vnější stěnu započteme $k$-krát. Pročež,
 abychom dostali skutečnou velikost $\# F_D$, musíme tuto $(k-1)$-krát odečíst.
 Z tohoto pozorování již snadno plyne, že
 \begin{equation}
  \label{eq:steny-komponent}
  \tag{$*$}
  \sum_{i=1}^{k} \# F_{D_i} = \# F_D + k - 1.
 \end{equation}
 Konečně, pro každou komponentu platí z předchozího vzorec
 \[
  \# V_i - \# E_i + \# F_{D_i} = 2.
 \]
 Sečtením přes všechny komponenty souvislosti dostaneme
 \[
  \sum_{i = 1}^{k} \# V_i - \sum_{i = 1}^{k} \# E_i + \sum_{i = 1}^{k} \#
  F_{D_i} = \sum_{i = 1}^{k} 2 = 2k.
 \]
 Dosazením \eqref{eq:steny-komponent} dostaneme
 \[
  \# V - \# E + \# F_D + k - 1 = 2k,
 \]
 z čehož pouhou úpravou
 \[
  \# V - \# E + \# F_D = k + 1,
 \]
 což zakončuje důkaz.
\end{proof}

\begin{figure}[h]
 \centering
 \begin{tikzpicture}[scale=1.5]
  \pgfdeclarelayer{background}
  \pgfdeclarelayer{foreground}
  \pgfsetlayers{background,main,foreground}

  \begin{pgfonlayer}{foreground}
   \node[vertex] (a) at (0,0) {};
   \node[vertex] (b) at (1,1.5) {};
   \node[vertex] (c) at (1,-1.5) {};
   \node[vertex] (d) at (2,0) {};
   \node[vertex] (e) at (1,0.5) {};
   \node[vertex] (f) at (1,-0.5) {};

   \draw[thick] (a) -- (b);
   \draw[thick] (a) -- (c);
   \draw[thick] (b) -- (d);
   \draw[thick] (c) -- (d);
   \draw[thick] (b) -- (e);
   \draw[thick] (c) -- (f);
  \end{pgfonlayer}
  \begin{pgfonlayer}{main}
   \draw[thick,->] (2.5,0) to node[midway,myred,yshift=3mm] {$+e$} (3.5,0);
  \end{pgfonlayer}
  \begin{pgfonlayer}{foreground}
   \node[vertex] (a1) at (4,0) {};
   \node[vertex] (b1) at (5,1.5) {};
   \node[vertex] (c1) at (5,-1.5) {};
   \node[vertex] (d1) at (6,0) {};
   \node[vertex] (e1) at (5,0.5) {};
   \node[vertex] (f1) at (5,-0.5) {};

   \draw[thick] (a1) -- (b1);
   \draw[thick] (a1) -- (c1);
   \draw[thick] (b1) -- (d1);
   \draw[thick] (c1) -- (d1);
   \draw[thick] (b1) -- (e1);
   \draw[thick] (c1) -- (f1);
   \draw[thick,myred] (e1) to node[midway,myred,xshift=2mm] {$e$} (f1);
  \end{pgfonlayer}
  \begin{pgfonlayer}{background}
   \fill[myblue,opacity=0.25] (a.center) -- (b.center) -- (d.center) --
    (c.center) -- (a.center);
   \node[myblue] at (1,0) {$\mathbf{f_1}$};
   \fill[myblue,opacity=0.25] (a1.center) -- (b1.center) -- (c1.center) --
    (a1.center);
   \fill[mygreen,opacity=0.25] (b1.center) -- (d1.center) -- (c1.center) --
    (b1.center);
   \node[myblue] at (4.5,0) {$\mathbf{f_1}$};
   \node[mygreen] at (5.5,0) {$\mathbf{f_1'}$};
  \end{pgfonlayer}
 \end{tikzpicture}

 \caption{Indukční krok v důkazu \myref{tvrzení}{claim:euleruv-vztah}.}
 \label{fig:eulerova-veta}
\end{figure}

\begin{remark}
 Důsledkem \hyperref[claim:euleruv-vztah]{Eulerova vztahu} je, že počet stěn
 rovinného grafu je \textbf{nezávislý na volbě nakreslení}. Přestože stěny
 samotné nemusejí být napříč nakresleními totožné, odpustíme sobě lehký
 prohřešek vůči rigorositě a budeme množinu stěn rovinného grafu $G$ značit
 obecně písmenem $F$, nedbajíce volby rovinného nakreslení. Máme-li totiž dvě
 rovinná nakreslení $D,D'$, pak mezi $F_D$ a $F_{D'}$ existuje bijekce. Prosté
 značení $F$ je tedy korektní \uv{až na bijekci}.
\end{remark}

Připomeňme problém tří domů a tří studní ze
\myref{sekce}{ssec:problem-tri-domu-a-tri-studni}. Na jejím závěru tvrdíme, že
nemá řešení, tedy, že neexistuje způsob, jak každý z tří domů spojit cestou s
každou ze tří studní tak, aby se cesty nekřížily. Nyní již tento problém umíme
matematicky zformulovat -- tvrdíme, že graf $K_{3,3}$ není rovinný.

Než toto tvrzení dokážeme, připomeneme rovněž, že zápisem $K_n$ myslíme úplný
graf na $n \in \N$ vrcholech, tj. graf, v němž je každý pár vrcholů spojen
hranou. Dále, zápis $K_{m,n}$ značí tzv. \emph{úplný bipartitní graf} na $m + n$
vrcholech. To je graf, jehož množinu vrcholů $V$ lze rozdělit na dvě disjunktní
podmnožiny $V_1$ a $V_2$ velikostí $m$ a $n$, respektive, takové, že mezi každým
párem vrcholů $v_1,v_2$, kde $v_1 \in V_1$ a $v_2 \in V_2$, vede hrana, ale mezi
dvěma vrcholy ze stejné množiny hrana nevede.

Snad nejdůležitější tvrzení teorie kreslení grafů,
\hyperref[thm:kuratowskiho]{Kuratowskiho věta}, tvrdí, že $K_5$ a $K_{3,3}$ jsou
\textbf{jediné dva nerovinné grafy} v tom smyslu, že každý nerovinný graf musí
jeden z nich obsahovat.

Tuto sekci zakončíme tím, že z \hyperref[claim:euleruv-vztah]{Eulerova vztahu}
odvodíme omezení na počet hran (v~závislosti na počtu vrcholů) rovinných grafů,
které $K_{3,3}$ ani $K_5$ nesplňují, což bude znamenat, že nejsou rovinné. Důkaz
opačné implikace, že každý nerovinný graf jeden z těchto grafů obsahuje je
mnohem těžší a je mu věnována \hyperref[sssec:kuratowskiho-veta]{následující
sekce}.

\begin{lemma}
 \label{lem:hrany-rovinnych-grafu}
 Ať $G = (V,E)$ je rovinný graf. Pak platí
 \begin{enumerate}[label=(\alph*)]
  \item $\# E \leq 3\# V - 6$;
  \item neobsahuje-li navíc $G$ trojúhelník (tj. $K_3$) jako podgraf, pak $\#
   E \leq 2\# V - 4$.
 \end{enumerate}
\end{lemma}
\begin{proof}
 Začneme důkazem části (a), část (b) se z něj pak snadnou modifikací odvodí.

 Nejprve uvážíme případ, kdy $G$ neobsahuje cyklus. Pak máme $\# F = 1$ a $\# E
 = \# V - k$ (buď z \hyperref[claim:euleruv-vztah]{Eulerova vztahu} nebo prostou
 úvahou). Nerovnost
 \[
  \# E = \# V - k \leq 3\# V - 6,
 \]
 platí právě tehdy, když $\# V \geq (6-k) / 2$. Jelikož $\# V$ i $k$ jsou
 přirozená čísla, tento vztah automaticky platí, kdykoli $\# V \geq 3$. To nám
 stačí, neboť všechny grafy na dvou a méně vrcholech jsou zřejmě rovinné.

 Nyní ať $G$ obsahuje cyklus. To zároveň znamená, že má aspoň dvě stěny. Vezměme
 libovolnou z nich a označme ji $f \in F$. Dále označme $E_f \subseteq E$
 množinu těch hran, které leží na hranici $f$ (tj. těch, které ji oddělují od
 ostatních stěn). Uvažme podgraf $G_f = (V,E_f)$. Ten je samozřejmě stále
 rovinný, neboť vznikl z $G$ odebráním nějakého počtu hran. Všimněme si dále, že
 $f$ je stále stěnou v $G_f$, bo hrany tvořící její hranici zůstaly nedotčeny, a
 tedy se nemohla $f$ zvětšit ani zmenšit. Protože nakreslení $G$ mělo aspoň dvě
 stěny, má nakreslení $G_f$ rovněž aspoň dvě stěny, a tudíž obsahuje cyklus.
 Tento cyklus je podmnožinou $E_f$, z čehož plyne, že $\# E_f \geq 3$, neboť
 cyklus má aspoň 3 hrany.

 Když nerovnost z předchozího odstavce sečteme přes všechny stěny $f \in F$,
 dostaneme
 \[
  \sum_{f \in F} \# E_f \geq \sum_{f \in F} 3 = 3 \# F.
 \]
 Zároveň, každá hrana je buď uvnitř stěny, nebo na její hranici, tj. náleží buď
 jedné nebo dvěma stěnám. Odtud máme další nerovnost,
 \[
  \sum_{f \in F} \# E_f \leq 2\# E.
 \]
 Sloučením obou nerovností obdržíme $3\# F \leq 2\# E$, neboli $\# F \leq 2\# E
 / 3$. Dosazením do \hyperref[claim:euleruv-vztah]{Eulerova vztahu} dostaneme
 \[
  \# E = \# V + \# F - k - 1 \leq \# V + \frac{2\# E}{3} - k - 1.
 \]
 Prostou úpravou pak
 \[
  \# E \leq 3\# V - 3k - 3.
 \]
 Protože $G$ má aspoň jednu komponentu ($k \geq 1$), plyne odtud
 \[
  \# E \leq 3\# V - 6,
 \]
 což dokazuje část (a). 

 Část (b) lze dokázat zcela obdobným způsobem. Stačí si pouze uvědomit, že
 neobsahuje-li $G$ trojúhelníky, pak každý cyklus má aspoň 4 hrany, a tudíž máme
 pro každou $f \in F$ nerovnost $\# E_f \geq 4$. Identická posloupnost úprav pak
 dá kýžený závěr.
\end{proof}

\begin{figure}[h]
 \centering
 \begin{tikzpicture}
  \pgfdeclarelayer{background}
  \pgfsetlayers{background,main}

  \node[vertex] (1) at (0,0) {};
  \node[vertex] (2) at (1,1) {};
  \node[vertex] (3) at (3,0) {};
  \node[vertex] (4) at (1,-1) {};
  \node[vertex] (5) at (2,0) {};
  \node[vertex] (6) at (2,1.5) {};
  \node[vertex] (7) at (3,1.2) {};

  \draw[ultra thick,myred] (1) -- (2);
  \draw[ultra thick,myred] (2) -- (3);
  \draw[ultra thick,myred] (3) -- (4);
  \draw[ultra thick,myred] (4) -- (1);
  \draw[thick] (3) -- (5);
  \draw[thick] (2) -- (5);
  \draw[thick] (2) -- (6);
  \draw[thick] (6) -- (7);

  \node[vertex] (11) at (-1,0) {};
  \node[vertex] (22) at (-1,2) {};
  \node[vertex] (33) at (2,3) {};
  \node[vertex] (44) at (3,3) {};
  \node[vertex] (55) at (4,1) {};
  \node[vertex] (66) at (5,0) {};
  \node[vertex] (77) at (3,-2) {};
  \node[vertex] (88) at (0.5,-2.5) {};
  
  \draw[ultra thick,myred] (11) -- (22);
  \draw[ultra thick,myred] (22) -- (33);
  \draw[ultra thick,myred] (33) -- (44);
  \draw[ultra thick,myred] (44) -- (55);
  \draw[ultra thick,myred] (55) -- (66);
  \draw[ultra thick,myred] (66) -- (77);
  \draw[ultra thick,myred] (77) -- (88);
  \draw[ultra thick,myred] (88) -- (11);

  \node[vertex] (r1) at (-2,3) {};
  \node[vertex] (r2) at (-2,4) {};
  \node[vertex] (r3) at (-3,3.5) {};
  
  \draw[thick] (22) -- (r1);
  \draw[thick] (r1) -- (r2);
  \draw[thick] (r2) -- (r3);
  \draw[thick] (r3) -- (r1);
  
  \node[vertex] (r4) at (-0.5,1.5) {};
  \draw[thick] (33) -- (r4);
  
  \node[vertex] (r5) at (4,3) {};
  \draw[thick] (44) -- (r5);
  \draw[thick] (r5) -- (55);
  
  \begin{pgfonlayer}{background}
  \fill[myblue,opacity=0.25] (11.center) -- (22.center) -- (33.center) --
   (44.center) -- (55.center) -- (66.center) -- (77.center) -- (88.center) --
   (11.center);
  \fill[white] (1.center) -- (2.center) -- (3.center) -- (4.center) --
   (1.center);
  \node[myblue] at (2.5,-1.2) {\Large $f$};
  \end{pgfonlayer}
 \end{tikzpicture}

 \caption{\clr{Množina $E_f$} v důkazu
 \myref{lemmatu}{lem:hrany-rovinnych-grafu}.}
 \label{fig:hrany-rovinnych-grafu}
\end{figure}

\begin{exercise}
 Najděte rovinný graf, jehož všechny vrcholy mají stupeň 5 (vizte
 \myref{definici}{def:stupen-vrcholu}).
\end{exercise}

\begin{exercise}
 Ukažte, že odhad $\# E \leq 2\# V - 4$ pro rovinný graf $G = (V,E)$ bez
 trojúhelníků je nejlepší možný. To jest, pro libovolné $n \in \N$ zkonstruujte
 rovinný graf s $n$ vrcholy a $2n - 4$ hranami.
\end{exercise}

\begin{exercise}
 Hra \uv{sprouts} vynalezená J. H. Conwayem a M. S. Patersonem má následující
 pravidla:
 \begin{itemize}
  \item Na papíře je nakresleno $n \in \N$ puntíků.
  \item Hráči se střídají po tazích, a kdo nemůže táhnout, prohrává.
  \item V každém tahu spojí hráč dva existující puntíky obloukem a někam na
   tento oblouk nakreslí nový puntík.
  \item Puntík je možné použít jako jeden z konců nového oblouku jedině tehdy,
   když v něm končí nanejvýš dva již nakreslené oblouky.
 \end{itemize}
 Na \myref{obrázku}{fig:priklad-sprouts} vizte příklad 4 prvních tahů pro $n =
 2$.
 \begin{figure}[H]
  \tikzset{
    vertex/.style = {shape=circle,fill,text=white,minimum size=6pt,inner
    sep=1pt}
  }
  \centering
  \begin{subfigure}[b]{.19\textwidth}
   \centering
   \begin{tikzpicture}[scale=0.5]
    \node[vertex] (a) at (180:1) {};
    \node[vertex] (b) at (0:1) {};
   \end{tikzpicture}
   \caption{Start.}
  \end{subfigure}
  \begin{subfigure}[b]{.19\textwidth}
   \centering
   \pgfdeclarelayer{background}
   \pgfsetlayers{background,main}
   \begin{tikzpicture}[scale=0.5]
    \node[vertex] (a) at (180:1) {};
    \node[vertex] (b) at (0:1) {};
    \node[vertex,myred] (c) at (90:1) {};

    \begin{pgfonlayer}{background}
     \draw[thick,myred] plot [smooth,tension=2] coordinates {(a) (c) (b)};
    \end{pgfonlayer}
   \end{tikzpicture}
   \caption{Tah 1.}
  \end{subfigure}
  \begin{subfigure}[b]{.19\textwidth}
   \centering
   \pgfdeclarelayer{background}
   \pgfsetlayers{background,main}
   \begin{tikzpicture}[scale=0.5]
    \node[vertex] (a) at (180:1) {};
    \node[vertex] (b) at (0:1) {};
    \node[vertex] (c) at (90:1) {};
    \node[vertex,myred] (d) at (-90:2) {};
    
    \begin{pgfonlayer}{background}
     \draw[thick] plot [smooth,tension=2] coordinates {(a) (c) (b)};
     \draw[thick,myred] plot [smooth,tension=2.5] coordinates {(a) (d) (b)};
    \end{pgfonlayer}
   \end{tikzpicture}
   \caption{Tah 2.}
  \end{subfigure}
  \begin{subfigure}[b]{.19\textwidth}
   \centering
   \pgfdeclarelayer{background}
   \pgfsetlayers{background,main}
   \begin{tikzpicture}[scale=0.5]
    \node[vertex] (a) at (180:1) {};
    \node[vertex] (b) at (0:1) {};
    \node[vertex] (c) at (90:1) {};
    \node[vertex] (d) at (-90:2) {};
    \node[vertex,myred] (e) at (0,0) {};
    
    \begin{pgfonlayer}{background}
     \draw[thick] plot [smooth,tension=2] coordinates {(a) (c) (b)};
     \draw[thick] plot [smooth,tension=2.5] coordinates {(a) (d) (b)};
     \draw[thick,myred] (a) -- (e);
     \draw[thick,myred] (e) -- (b);
    \end{pgfonlayer}
   \end{tikzpicture}
   \caption{Tah 3.}
  \end{subfigure}
  \begin{subfigure}[b]{.19\textwidth}
   \centering
   \pgfdeclarelayer{background}
   \pgfsetlayers{background,main}
   \begin{tikzpicture}[scale=0.5]
    \node[vertex] (a) at (180:1) {};
    \node[vertex] (b) at (0:1) {};
    \node[vertex] (c) at (90:1) {};
    \node[vertex] (d) at (-90:2) {};
    \node[vertex] (e) at (0,0) {};
    \node[vertex,myred] (f) at (2,-0.5) {};
 
    \begin{pgfonlayer}{background}
     \draw[thick] plot [smooth,tension=2] coordinates {(a) (c) (b)};
     \draw[thick] plot [smooth,tension=2.5] coordinates {(a) (d) (b)};
     \draw[thick] (a) -- (e);
     \draw[thick] (e) -- (b);
     \draw[thick,myred] (c) arc (106.35:-106.35:1.562);
    \end{pgfonlayer}
   \end{tikzpicture}
   \caption{Tah 3.}
  \end{subfigure}
  \caption{Příklad prvních čtyř tahů hry \emph{sprouts} pro $n = 2$.}
  \label{fig:priklad-sprouts}
 \end{figure}
 Dokažte, že
 \begin{enumerate}[label=(\alph*)]
  \item pro $n$ počátečních puntíků má tato hra maximálně $3n - 1$ tahů;
  \item (těžké) pro $n$ počátečních puntíků má tato hra \textbf{nejméně} $2n$
   tahů (bez ohledu na strategii hráčů).
 \end{enumerate}
\end{exercise}
>>>>>>> Stashed changes
